\documentclass[a4paper,12pt]{article}
\usepackage{graphicx}
\usepackage{float}
\usepackage{booktabs}
\usepackage{amsmath}
\usepackage{geometry}
\usepackage{caption}
\usepackage{array}
\usepackage{hyperref}
\usepackage{multirow}
\usepackage{longtable}
\usepackage{tabularx}
\usepackage{booktabs}

\geometry{margin=1in}

\begin{document}

\begin{center}
     \textbf{Sri Sivasubramaniya Nadar College of Engineering, Chennai}\\
    (An Autonomous Institution affiliated to Anna University)\\[4pt]
     Department of Computer Science and Engineering\\[8pt]
    \textbf{CAT Assignment 1}\\[10pt]
\end{center}
    \begin{tabular}{rl}
        \textbf{Degree \& Branch:} & M.Tech (Integrated) Computer Science and Engineering \\
        \textbf{Semester:} & V \\
        \textbf{Subject Code \& Name:} & ICS1502 - Introduction to Machine Learning \\
        \textbf{Academic Year:} & 2025–2026 (Odd Semester) \\
        \textbf{Batch:} & 2023–2028 \\
        \textbf{Name:} & Jayasree R \\
        \textbf{Register Number:} & 3122 23 7001 019 \\
    \end{tabular}

\vspace{0.5cm}
\hrule
\vspace{0.5cm}

\section*{1. Programming Assignment}

\subsection*{1.1 Regression: Mobile Phone Price Prediction}
\textbf{Aim:} To implement and analyze a Linear Regression model for predicting mobile phone prices using both closed-form and gradient descent methods, with and without L2 regularization, and to evaluate their performance under various preprocessing settings.

\subsubsection*{Objectives}
\begin{itemize}
    \item To construct data and label matrices for linear regression.
    \item To implement regression using closed-form and gradient descent solutions.
    \item To study the effect of regularization (L2 penalty) on model generalization.
    \item To compare model performance before and after data standardization.
    \item To interpret model weights and understand feature importance.
\end{itemize}

\subsubsection*{Mathematical Description}
Given a dataset $(X, y)$ with $m$ samples and $n$ features, the linear regression model predicts:
\[
\hat{y} = X\theta
\]
The closed-form solution (Ordinary Least Squares) is:
\[
\theta = (X^T X)^{-1} X^T y
\]
For Ridge Regression (L2 regularization):
\[
\theta = (X^T X + \lambda I)^{-1} X^T y
\]
For Gradient Descent:
\[
\theta := \theta - \alpha \frac{1}{m} X^T (X\theta - y)
\]
where $\alpha$ is the learning rate and $\lambda$ controls the regularization strength.

\subsubsection*{Results}

\textbf{Exploratory Data Analysis (EDA):}
\begin{itemize}
    \item Checked for missing values and feature distributions.
    \item Generated correlation heatmap to identify feature relationships.
\end{itemize}

\textbf{Pairplot and Correlation Heatmap:}

\begin{figure}[H]
    \centering
    \includegraphics[width=0.8\linewidth]{theory_ass1/screenshots/pairplot_regression.png}
    \caption{Pairplot of features in Mobile Price dataset}
\end{figure}

\begin{figure}[H]
    \centering
    \includegraphics[width=0.8\linewidth]{theory_ass1/screenshots/correlation_heatmap_regression.png}
    \caption{Feature correlation heatmap}
\end{figure}

\textbf{Performance Metrics:}

\begin{itemize}
    \item \textbf{Closed-form Solution:}\\
    MSE = 23062.23, RMSE = 151.86, MAE = 130.32, $R^2$ = 0.9593
    \item \textbf{Gradient Descent:}\\
    MSE = 8.2462e+10, RMSE = 287162.01, MAE = 264011.00, $R^2$ = -145463.13
    \item \textbf{Ridge (Closed-form, $\lambda$=10):}\\
    Test RMSE = 161.09
    \item \textbf{Ridge (Gradient Descent, $\lambda$=10):}\\
    Test RMSE = 1909.43
\end{itemize}

\textbf{Comparison Across $\lambda$ Values:}

\begin{longtable}{cccccc}
\toprule
\textbf{Lambda} & \textbf{No-Scale RMSE} & \textbf{Scale RMSE} & \textbf{No-Scale $R^2$} & \textbf{Scale $R^2$} \\
\midrule
0.00 & 151.86 & 151.86 & 0.9593 & 0.9593 \\
0.01 & 151.86 & 151.87 & 0.9593 & 0.9593 \\
0.10 & 151.87 & 151.92 & 0.9593 & 0.9593 \\
1.00 & 152.01 & 152.49 & 0.9592 & 0.9589 \\
10.00 & 154.24 & 161.09 & 0.9580 & 0.9542 \\
100.00 & 166.65 & 230.26 & 0.9510 & 0.9065 \\
\bottomrule
\end{longtable}

\textbf{Predicted vs Actual Plots:}
\begin{figure}[H]
    \centering
    \includegraphics[width=1\linewidth]{theory_ass1/screenshots/actual_vs_predicted.png}
    \caption{Predicted vs Actual Plot }
\end{figure}

\begin{figure}[H]
    \centering
    \includegraphics[width=0.9\linewidth]{theory_ass1/screenshots/lamba.png}
    \caption{Predicted vs Actual for Ridge Regression with different $\lambda$ values}
\end{figure}

\begin{figure}[H]
    \centering
    \includegraphics[width=0.7\linewidth]{theory_ass1/screenshots/features.png}
    \caption{Feature Importance Plot}
\end{figure}


\subsubsection*{Inference}
\begin{itemize}
    \item The closed-form linear regression performed well with a high $R^2$ (0.9593), showing strong predictive power.
    \item Gradient Descent failed to converge due to improper learning rate or lack of normalization.
    \item Ridge regularization slightly reduced RMSE, preventing overfitting for larger $\lambda$ values.
    \item Excessively large $\lambda$ degraded model performance, emphasizing the need for regularization tuning.
    \item Feature importance analysis showed that certain features had strong linear relationships with price, consistent with correlation heatmap findings.
    \item Standardization of data improved the stability of results for smaller $\lambda$ values.
\end{itemize}


\newpage
\subsection*{1.2 Linear Classification: Bank Note Authentication}

\textbf{Aim:} To fit a linear classification model for banknote authentication using logistic regression and neural networks, evaluate L2 regularization, and analyze model robustness under outlier injection.

\subsubsection*{Objectives}
\begin{itemize}
    \item To train and test a linear classifier on the Bank Note Authentication dataset.
    \item To compare model accuracy with and without L2 regularization.
    \item To analyze the effect of outliers on model performance.
    \item To visualize 3D decision boundaries using key features.
\end{itemize}

\subsubsection*{Mathematical Description}
The logistic regression model predicts the probability:
\[
P(y=1|x) = \sigma(w^T x + b)
\]
where $\sigma(z) = \frac{1}{1 + e^{-z}}$.

The cost function with L2 regularization:
\[
J(w) = -\frac{1}{m} \sum_{i=1}^{m} [y^{(i)} \log(\hat{y}^{(i)}) + (1 - y^{(i)}) \log(1 - \hat{y}^{(i)})] + \frac{\lambda}{2m}\|w\|^2
\]

\subsubsection*{Results}

\textbf{Exploratory Data Analysis:}
\begin{itemize}
    \item Dataset was clean and well-balanced across classes.
    \item Features displayed strong linear separability, as seen in pairplots.
\end{itemize}

\textbf{Pairplot and Heatmap:}
\begin{figure}[H]
    \centering
    \includegraphics[width=0.8\linewidth]{theory_ass1/screenshots/pairplot_classification.png}
    \caption{Feature pairplot for Bank Note Authentication dataset}
\end{figure}

\begin{figure}[H]
    \centering
    \includegraphics[width=0.8\linewidth]{theory_ass1/screenshots/correlation_classification.png}
    \caption{Correlation heatmap of features}
\end{figure}

\textbf{Model Performance:}


\begin{table}[H]
\centering
\caption{Performance Comparison of Models}
\begin{tabular}{lcccc}
\toprule
\textbf{Model} & \textbf{Train Acc} & \textbf{Test Acc} & \textbf{Precision / Recall / F1}  \\
\midrule
No Regularization & 0.9909 & 0.9854 & 0.98–0.99  \\
L2 (C=1.0)        & 0.9845 & 0.9709 & 0.97       \\
MLP (1 layer)      & 1.0000 & 1.0000 & 1.00     \\
\bottomrule
\end{tabular}
\label{tab:model_performance}
\end{table}



\textbf{Feature Importance Ranking:}

\begin{longtable}{lc}
\toprule
\textbf{Feature} & \textbf{Importance Score} \\
\midrule
Skewness & 4.8089 \\
Variance & 4.6579 \\
Curtosis & 4.3977 \\
Entropy & 0.2552 \\
\bottomrule
\end{longtable}

\textbf{Effect of Outliers:}

\begin{table}[H]
\centering
\begin{tabular}{lcc}
\toprule
\textbf{Model} & \textbf{Train Accuracy} & \textbf{Test Accuracy} \\
\midrule
Clean Model & 0.9845 & 0.9709 \\
Outlier Model & 0.8441 & 0.8182 \\
\bottomrule
\end{tabular}
\caption{Impact of outlier injection on model performance}
\end{table}

\textbf{Confusion Matrices:}

\begin{table}[H]
\centering
\caption{Comparison of Classification Metrics between Clean and Outlier Models}
\begin{tabular}{lcccc}
\hline
\textbf{Metric} & \textbf{Class} & \textbf{Clean Model} & \textbf{Outlier Model} & \textbf{Difference ($\Delta$)} \\ 
\hline
\textbf{Precision} & 0.0 & 1.00 & 0.86 & -0.14 \\ 
                   & 1.0 & 0.94 & 0.78 & -0.16 \\ 
\hline
\textbf{Recall}    & 0.0 & 0.95 & 0.81 & -0.14 \\ 
                   & 1.0 & 1.00 & 0.83 & -0.17 \\ 
\hline
\textbf{F1-Score}  & 0.0 & 0.97 & 0.83 & -0.14 \\ 
                   & 1.0 & 0.97 & 0.80 & -0.17 \\ 
\hline
\textbf{Accuracy} & -- & 0.97 & 0.82 & -0.15 \\ 
\textbf{Macro Avg} & -- & 0.97 & 0.82 & -0.15 \\ 
\textbf{Weighted Avg} & -- & 0.97 & 0.82 & -0.15 \\ 
\hline
\end{tabular}
\label{tab:classification_metrics_comparison}
\end{table}



\textbf{3D Visualization of Decision Boundaries:}
\begin{figure}[H]
    \centering
    \includegraphics[width=0.7\linewidth]{theory_ass1/screenshots/3d_plot_classification.png}
    \caption{3D visualization using top three important features}
\end{figure}

\subsubsection*{Inference}
\begin{itemize}
    \item The dataset was linearly separable; hence, linear models achieved high accuracy.
    \item Without regularization, the classifier achieved 98.5\% test accuracy, showing minimal overfitting.
    \item With L2 regularization, accuracy dropped slightly but model stability improved.
    \item The MLP model achieved perfect accuracy, confirming non-linear models’ superior fitting capacity.
    \item Feature ranking indicated \textit{skewness}, \textit{variance}, and \textit{curtosis} as the most discriminative attributes.
    \item Injecting outliers drastically reduced accuracy (from 0.97 to 0.82), revealing that linear models are sensitive to anomalous data.
    \item The decision boundary visualizations confirmed that most samples were well-separated in 3D space.
\end{itemize}

\section*{Learning Outcomes}
\begin{itemize}
    \item Understood matrix-based implementation of linear regression.
    \item Gained practical insight into gradient descent convergence behavior.
    \item Learned to tune and interpret regularization effects on bias–variance trade-off.
    \item Acquired experience in evaluating linear classifiers using performance metrics and confusion matrices.
    \item Observed the impact of outliers and importance of data quality in supervised learning.
    \item Enhanced understanding of data visualization and feature interpretation through pairplots, heatmaps, and 3D projections.
\end{itemize}
\section*{GitHub Repository}

The complete source code for this assignment is available in the following repository: 
\href{https://github.com/R-Jayasree/Machine-Learning} \textbf{https://github.com/R-Jayasree/Machine-Learning}

\end{document}
